This guide explains how LaTeX works and how to use it well, starting from nothing. It\textquotesingle s also my own reference, so that I can rebuild my projects from scratch if I ever forget how they work.

LaTeX suits long documents with maths, numbered sections and references, because you describe the document\textquotesingle s structure and it handles the layout. You write plain text, which also means every version can be tracked and compared.

\section{How to use it}\label{how-to-use-it}

Read in order if you\textquotesingle re new. Each chapter builds on the last, and the code shown compiles as written. The \href{examples/latex-examples.zip}{example files} include all of it, and the complete report comes together in the chapter on \href{references/}{references}.

Prefer one long page? \href{full/}{Read the whole guide on one page}.

\section{Getting LaTeX}\label{getting-latex}

\begin{itemize}
\tightlist
\item
  \textbf{To try it with nothing installed:} \href{https://www.overleaf.com}{Overleaf}\footnote{\texttt{https://www.overleaf.com}} runs in the browser.
\item
  \textbf{On your computer:} \href{https://miktex.org}{MiKTeX}\footnote{\texttt{https://miktex.org}} on Windows, \href{https://www.tug.org/mactex/}{MacTeX}\footnote{\texttt{https://www.tug.org/mactex/}} on macOS, or \href{https://www.tug.org/texlive/}{TeX Live}\footnote{\texttt{https://www.tug.org/texlive/}} on Linux.
\item
  \textbf{An editor:} VS Code with the LaTeX Workshop extension, or TeXstudio.
\end{itemize}

The \href{https://www.latex-project.org/help/documentation/}{official LaTeX documentation}\footnote{\texttt{https://www.latex-project.org/help/documentation/}} and the \href{https://en.wikibooks.org/wiki/LaTeX}{LaTeX Wikibook}\footnote{\texttt{https://en.wikibooks.org/wiki/LaTeX}} are good references to keep at hand.

